\documentclass[12pt,a4paper]{article}
\usepackage{ctex}
\usepackage{amsmath,amssymb}
\usepackage{booktabs}
\title{基于深度学习的图像分类研究}
\author{张三}
\date{\today}

\begin{document}

\maketitle

\section{引言}
本节介绍研究背景。行内公式 $E = mc^2$ 用于说明质能关系，
\textbf{加粗文本}与\textit{斜体文本}也应正常显示。

行间公式如下：
\[
\int_{0}^{1} x^2 \, dx = \frac{1}{3}
\]

\subsection{相关工作}
现有方法主要有三类：

\begin{itemize}
  \item 基于卷积神经网络的方法
  \item 基于\texttt{Transformer}的方法
  \item 混合架构方法
\end{itemize}

\subsection{实验设置}

\begin{enumerate}
  \item 数据集：CIFAR-10
  \item 优化器：Adam，学习率 0.001
\end{enumerate}

\section{实验结果}
% 这是一条注释，不应出现在 Word 里

\begin{table}[htbp]
\centering
\begin{tabular}{lcc}
\toprule
方法 & 准确率 & 参数量 \\
\midrule
ResNet-18 & 95.2\% & 11M \\
ViT-B & 97.8\% & 86M \\
\bottomrule
\end{tabular}
\caption{不同方法对比}
\label{tab:cmp}
\end{table}

结果如表 \ref{tab:cmp} 所示，其中 50\% \& \$100 这类特殊字符应被正确还原。

\section{结论}
本文提出的方法在准确率上提升了 2.6\%。

\end{document}
